
%%%%%%%%%%%%%%%%%%%%%%%%%%%%%%%%%%%%%%%%
%           main body
%%%%%%%%%%%%%%%%%%%%%%%%%%%%%%%%%%%%%%%%

\usepackage{tabu}

%-----------------------------------
%	TITLE PAGE
%-----------------------------------

\title[Probability \& Statistics]{\LARGE \bfseries 第三部分\quad 概率与统计} % The short title appears at the bottom of every slide, the full title is only on the title page
\author[Hypothesis Test] % Your institution as it will appear on the bottom of every slide, may be shorthand to save space
{\Large \bfseries \S 9. 假设检验}
% \author{Singular Value Decomposition} % Your name
% \date{\today} % Date, can be changed to a custom date
\date{}

%------------------------------------------------

\begin{document}


%------------------------------------------------

\begin{frame}

\titlepage % Print the title page as the first slide

\end{frame}

%------------------------------------------------

\begin{frame}

\end{frame}

%------------------------------------------------

\begin{frame}
\frametitle{内容提要} % Table of contents slide, comment this block out to remove it
\tableofcontents % Throughout your presentation, if you choose to use \section{} and \subsection{} commands, these will automatically be printed on this slide as an overview of your presentation
\end{frame}

%-----------------------------------
%	PRESENTATION SLIDES
%-----------------------------------

\section{矩估计法}

%------------------------------------------------

\begin{frame}

\begin{block}{矩估计}

\end{block}

\end{frame}

%------------------------------------------------

\begin{frame}

\begin{block}{矩估计}

\end{block}

\end{frame}

%------------------------------------------------

\begin{frame}

\begin{block}{矩估计例子}

\end{block}

\end{frame}

%------------------------------------------------

\begin{frame}

\begin{block}{矩估计例子（续）}

\end{block}

\end{frame}


%------------------------------------------------

\section{最大似然估计法}

%------------------------------------------------

\begin{frame}

最大似然估计法的思想是：利用已得到的抽样结果，寻找使这个结果出现的可能性最大的参数的值作为参数的估计。

\begin{block}{似然函数，Likelihood Function}

\end{block}

\end{frame}

%------------------------------------------------

\begin{frame}

\begin{block}{最大似然估计法，Maximum Likelihood Estimation, MLE}

\end{block}

\end{frame}

%------------------------------------------------

\begin{frame}

\begin{block}{似然函数求极大值}

\end{block}

\end{frame}

%------------------------------------------------

\begin{frame}

\begin{block}{最大似然估计法例子}

\end{block}

\end{frame}

%------------------------------------------------

\begin{frame}

\begin{block}{最大似然估计法例子（续）}

\end{block}

\end{frame}

%------------------------------------------------

\begin{frame}

\begin{block}{最小二乘法}

\end{block}

\end{frame}

%------------------------------------------------

\begin{frame}

\begin{block}{最小二乘法}

\end{block}

\end{frame}

%------------------------------------------------

% \begin{frame}

% \begin{block}{Vandermonde行列式的计算}
% \begin{enumerate}
% % \addtocounter{enumi}{0}
% \item 把第$n-1$行乘以$-a_1$加到第$n$行，再把第$n-2$行乘以$-a_1$加到第$n-1$行。按照这种顺序依次向上，直到用第$2$行乘以$-a_1$加到第$1$行。由行列式性质，$V_n$的值不变。
% \[
% V_n = \det \begin{pmatrix}
% 1 & 1 & \cdots & 1 \\
% 0 & a_2-a_1 & \cdots & a_n-a_1 \\
% \vdots & \vdots & & \vdots \\
% 0 & a_2^{n-2}(a_2-a_1) & \cdots & a_n^{n-2}(a_n-a_1) \\
% \end{pmatrix}
% \]
% \end{enumerate}
% \end{block}

% \end{frame}

%------------------------------------------------
% % closing page
% \begin{frame}
% \HUGE{\centerline{\bfseries {\color{gray} The} {\color{zkblue} End}}}

% \pause

% \vspace{1.2em}

% \Large{\centerline{\bfseries {\color{gray}Thank} \quad {\color{zkblue} You}}}

% \end{frame}

%----------------------------------------------------------------------------------------

\end{document}
